\chapter{Marco Teórico}
\label{chap:theo}

\section{Marco Conceptual}
\label{sec:conceptos}

\subsection{IDS en redes IoT/IIoT}
Un sistema de detección de intrusiones de red (NIDS) monitorea tráfico y eventos para identificar comportamientos maliciosos o anómalos. En IoT/IIoT, esta tarea es más exigente que en redes tradicionales porque los dispositivos tienen recursos limitados, usan protocolos livianos, presentan patrones de tráfico heterogéneos y suelen estar distribuidos en dominios administrativos distintos. Por ello, la detección debe evaluarse no solo por \textit{accuracy}, sino también por latencia, memoria, comunicación y estabilidad operacional \cite{Tekin2023_OnDeviceEnergy, Amuthadevi2025}. Surveys recientes sobre seguridad IIoT habilitada por ML insisten en que la diversidad de capas (sensor, gateway, fog, cloud) impide soluciones monolíticas y favorece arquitecturas jerárquicas con responsabilidades repartidas \cite{Ni2024_IIoT_Security_Survey, Amanullah2020_DL_BigData_IoT_Security}. Frameworks anteriores combinaron SDN/NFV con ML para coordinar respuesta y monitoreo (\emph{e.g.,}\cite{Bagaa2020_ML_IoT_Framework, Huda2018_DBN_SCADA}), e incluso frameworks ML basados en cloud para gestionar la seguridad de IoT~\cite{Mohamed2024_ML_CloudIoT}, pero rara vez mostraron despliegue sobre microcontroladores reales.

\subsection{Representación por flujos}
El tráfico de red puede representarse como paquetes, capturas PCAP, logs o flujos. Esta tesis utiliza flujos unidireccionales porque son compactos, tabulares y computables en el borde. Cada flujo se transforma en 13 características numéricas relacionadas con conteo de paquetes, tiempos entre paquetes, tamaños de paquetes, bytes totales y banderas TCP. Esta representación sacrifica parte del detalle secuencial de un PCAP completo, pero permite inferencia TinyML en ESP32-S3 y reduce el volumen de comunicación.

\subsection{TinyML y modelos compactos}
TinyML agrupa técnicas para ejecutar modelos de aprendizaje automático en hardware restringido. En el contexto de esta tesis, el criterio de diseño es que el modelo pueda ejecutarse directamente en microcontroladores ESP32-S3 sin depender de aceleradores externos ni transferir los datos crudos al servidor. La revisión de Dini~\textit{et~al.}~\cite{Dini2024_EmbeddedIIoT_Review} resume el ecosistema actual de \emph{toolchains} (TensorFlow Lite Micro, STM32Cube.AI, Vitis~AI, etc.) y muestra que la integración exitosa depende tanto de la arquitectura del modelo como del soporte HW/SW del MCU. Ficco~\textit{et~al.}~\cite{Ficco2024_FL_TinyML_Onboard} demostraron que el entrenamiento on-board en dispositivos como ESP32 y Arduino MKR1010 es viable cuando se combina FL con \emph{transfer learning}, alcanzando accuracy comparable al entrenamiento centralizado (86.48\% en clasificación) con latencia y consumo controlados. Los modelos evaluados en esta tesis incluyen MLP, CNN-1D, Residual-MLP y Tiny Transformer; sin embargo, el despliegue principal utiliza un MLP compacto por su equilibrio entre tamaño, precisión y simplicidad de inferencia, en línea con las recomendaciones de la literatura para clasificación tabular en MCUs \cite{HeydariMahmoud2025_TinyML_IDS, Dini2024_EmbeddedIIoT_Review}.

\subsection{Aprendizaje federado jerárquico}
El aprendizaje federado (FL) permite entrenar modelos de forma distribuida compartiendo actualizaciones en lugar de datos crudos. En escenarios IoT, una arquitectura jerárquica (HFL) resulta natural porque los nodos edge pueden comunicarse con gateways cercanos y estos, a su vez, con un coordinador central. Este diseño reduce exposición de datos, organiza la comunicación por capas y permite agregar actualizaciones localmente antes de enviarlas al servidor global \cite{McMahan2017FedAvg, Imteaj2022_FL_ResourceConstrainedSurvey}. Hajj~\textit{et~al.}~\cite{Hajj2023_CrossLayerFL_IDS} muestran que un IDS federado \emph{cross-layer} para IoT puede alcanzar precisión competitiva con menor consumo de comunicación, y Chaurasia~\textit{et~al.}~\cite{Chaurasia2024_FL_ResNet_IIoT} evidencian beneficios al usar redes residuales como base del cliente FL en IIoT.

La literatura advierte que FL en IoT enfrenta datos \textit{non-IID}, conectividad irregular y heterogeneidad de hardware \cite{Albanbay2025, Ramadan2025}. En esta tesis, esas condiciones se reproducen parcialmente mediante nodos con distribuciones distintas: un nodo genera tráfico normal y otro genera una mezcla de tráfico normal, fuerza bruta MQTT y escaneo. Adicionalmente, Tran~\textit{et~al.}~\cite{Tran2026_EdgeTrustShard} cuantifican una brecha de memoria de hasta tres órdenes de magnitud entre los métodos de agregación bizantino-resilientes (512\,MB--2\,GB) y los MCU comerciales ($\sim$264\,KB de SRAM), evidenciando la necesidad de descomponer el problema en capas. La arquitectura de esta tesis adopta esa filosofía sin requerir blockchain ni consenso PoP, dejando la robustez bizantina como trabajo futuro.

\subsection{Cifrado ligero y ASCON-128}
Aunque FL reduce la centralización de datos, los pesos del modelo y mensajes de control siguen siendo sensibles. Un atacante que intercepte o modifique actualizaciones podría inferir información, degradar el modelo o introducir comportamiento malicioso. Por esta razón, la arquitectura integra ASCON-128 como cifrado autenticado ligero para proteger confidencialidad e integridad en los canales MQTT y HTTP del ciclo federado.

ASCON fue seleccionado por NIST en febrero de 2023 como estándar de criptografía ligera, después de un proceso multi-anual frente a candidatos como GIFT-COFB, Sparkle, Romulus y TinyJambu \cite{NISTSP8002322025, Kaur2025_ASCON_Survey}. Internamente, ASCON aplica una permutación de 320 bits estructurada en cinco palabras de 64 bits, con dos variantes principales: ASCON-128 (rate 64 bits) y ASCON-128a (rate 128 bits). El survey de Kaur~\textit{et~al.}~\cite{Kaur2025_ASCON_Survey} resume implementaciones FPGA/ASIC, ataques diferenciales, ataques de canal lateral y contramedidas, ofreciendo una guía consolidada para diseñadores de sistemas embebidos. Su diseño resulta apropiado para mensajes pequeños como features, pesos parciales y modelos globales \cite{NISTSP8002322025, Gusita2025_LWC_Survey}.

Una alternativa frecuente en la industria son los protocolos \emph{ultra-lightweight} basados en permutaciones \emph{ad-hoc} (XOR, rotaciones, shifts), pero Mirzaali~\textit{et~al.}~\cite{Mirzaali2025_LWC_Permutations} demuestran experimentalmente sobre tres Raspberry Pi que estos esquemas suelen ser vulnerables a ataques de divulgación completa de secretos y desincronización, incluso cuando la estructura del protocolo parece sólida. Esto refuerza la decisión de usar un AEAD estandarizado en lugar de criptografía hecha en casa. En esta tesis se mide explícitamente su overhead frente a una versión funcionalmente equivalente sin cifrado, replicando el patrón experimental de varios trabajos previos sobre Edge--Cloud con cifrado por atributos \cite{SaravanaBalaji2023_EdgeIoT_DL}.

\subsection{Agregación FOG}
El modo FOG extiende la arquitectura estándar con comunicación entre gateways Raspberry Pi. Un gateway líder recibe actualizaciones de gateways pares mediante MQTT, realiza un FedAvg intermedio y envía al servidor central una actualización preagregada. Este esquema aproxima escenarios reales con múltiples zonas o sedes, donde una capa fog reduce tráfico hacia la nube y permite coordinación local antes de la agregación global. Wang~\textit{et~al.}~\cite{Wang2019_DataFusion_CPSS} discuten cómo la fusión de datos en sistemas ciber-físico-sociales se beneficia de pre-agregación cercana a la fuente, y EdgeTrust-Shard~\cite{Tran2026_EdgeTrustShard} demuestra que distribuir complejidad computacional por la topología es indispensable cuando los participantes son MCUs. El presente trabajo opera en una jerarquía de tres niveles (PC, Pi, ESP32) que materializa esa idea sin introducir blockchain.

\section{Estado del Arte}
\label{sec:antecedentes}

La literatura revisada se agrupa en cuatro líneas, articuladas alrededor de las cuatro dimensiones del sistema (Edge, Fog/FL, criptografía y datasets/etiquetado).

\subsection{IDS y TinyML para IoT/IIoT}
La primera línea aborda IDS para IoT/IIoT y destaca que los modelos de alta precisión no siempre son desplegables en el borde. Trabajos recientes sobre TinyML e IDS enfatizan que la memoria, la latencia y el costo energético deben reportarse junto con las métricas de clasificación \cite{Tekin2023_OnDeviceEnergy, Alwaisi2024, Amuthadevi2025, HeydariMahmoud2025_TinyML_IDS}. Surveys más generales sobre IIoT, DL y big data en seguridad IoT \cite{Amanullah2020_DL_BigData_IoT_Security, Ni2024_IIoT_Security_Survey} muestran que la mayoría de los pipelines reportados se apoyan en servidores capaces y carecen de despliegue real en MCUs. La revisión de Dini~\textit{et~al.}~\cite{Dini2024_EmbeddedIIoT_Review} compila explícitamente la oferta de \emph{toolchains} TinyML (TFLM, STM32Cube.AI, Vitis~AI) y los compromisos al integrarlos en plataformas embebidas.

\subsection{Aprendizaje federado para seguridad IoT}
La segunda línea corresponde al aprendizaje federado aplicado a seguridad IoT. FL permite entrenar modelos sin centralizar datos, pero introduce retos de comunicación, datos \textit{non-IID} y participación parcial. Estudios recientes sobre FL para IDS muestran que la precisión puede degradarse frente al entrenamiento centralizado, pero que la reducción de datos compartidos puede justificar el intercambio en escenarios con restricciones de privacidad \cite{Imteaj2022_FL_ResourceConstrainedSurvey, Albanbay2025, Hajj2023_CrossLayerFL_IDS, Roy2023_FL_AdaptedModel}. Chaurasia~\textit{et~al.}~\cite{Chaurasia2024_FL_ResNet_IIoT} usan redes residuales como cliente FL en IIoT, mientras que Ficco~\textit{et~al.}~\cite{Ficco2024_FL_TinyML_Onboard} validan FL+TL para entrenamiento on-board en MCUs y Raspberry Pi~3, comparando contra TFLite. EdgeTrust-Shard~\cite{Tran2026_EdgeTrustShard} aporta una visión más reciente: cuando los participantes son MCUs, las defensas Byzantine-robust convencionales (Krum, FLTrust, Fed-NGA) son inalcanzables sin descomponer el problema en una jerarquía como la propuesta en este trabajo.

\subsection{Criptografía ligera y autenticación}
La tercera línea es criptografía ligera. ASCON y otros esquemas de bajo costo buscan proteger dispositivos restringidos sin introducir overhead prohibitivo. Para un IDS federado, esta línea es especialmente relevante porque las actualizaciones de modelo son parte del activo de seguridad y deben protegerse durante el tránsito \cite{NISTSP8002322025, Gusita2025_LWC_Survey, Kaur2025_ASCON_Survey}. El survey de Kaur~\textit{et~al.}~\cite{Kaur2025_ASCON_Survey} provee una taxonomía de implementaciones FPGA/ASIC, ataques (diferenciales, cube, fault injection, side-channel) y contramedidas, lo que orienta decisiones de diseño y trabajos futuros. En contraste, Mirzaali~\textit{et~al.}~\cite{Mirzaali2025_LWC_Permutations} muestran cómo permutaciones ultra-ligeras \emph{ad-hoc} para autenticación RFID/IoT se rompen experimentalmente en hardware (RPi + Wi-Fi). Mecanismos complementarios de control de acceso \cite{Amoon2020_RRAC_IoT} y cifrado por atributos \cite{SaravanaBalaji2023_EdgeIoT_DL} se han propuesto para escenarios sanitarios y de smart cities, pero ninguno integra cifrado y FL sobre la misma plataforma.

\subsection{Datasets y representación del tráfico}
La cuarta línea cubre datasets y representación de tráfico. UNSW-NB15, Bot-IoT, TON\_IoT, Edge-IIoTset, IoT-23 y CICIoT2023 \cite{Moustafa2020_TONIoT, UNSW_UNSWNB15_Download, UNSW_BotIoT_Download, UNSW_TONIoT_Download, Ferrag2022_EdgeIIoTset, Stratosphere_IoT23, Neto2023_CICIoT2023} ofrecen capturas etiquetadas con diversos protocolos y vectores de ataque. Trabajos recientes \cite{Rani2024_IoT_IDS_Datasets, Raminenei2023_ML_IoT_UNSWNB15} proponen consolidar/integrar varios datasets para evitar sobreajuste a una sola fuente. En esta tesis se trabaja con tres datasets de flujos unidireccionales (\texttt{normal}, \texttt{mqtt\_bruteforce}, \texttt{scan\_A}) representativos del tráfico que un nodo MQTT puede observar en un laboratorio o smart home.

\section{Síntesis de la Brecha}

Los trabajos existentes suelen cubrir por separado TinyML, FL o cifrado ligero. La brecha abordada en esta tesis es la integración y medición de un ciclo completo Edge--Fog--Cloud sobre hardware real, con inferencia en ESP32-S3, entrenamiento local en Raspberry Pi, agregación global en PC, protección ASCON-128, línea base sin cifrado, variante CNN y modo FOG. Esta evaluación conjunta permite discutir el sistema como una arquitectura operativa y no solo como un clasificador aislado. Cuatro contribuciones específicas distinguen este trabajo:
\begin{enumerate}
    \item Implementación interoperable de ASCON-128 en C++ (ESP32) y Python (RPi/PC), instrumentada con métricas de tiempo y tamaño por canal y por ronda.
    \item Comparación experimental directa, sobre el mismo hardware y mismas condiciones, entre la rama con ASCON y la rama sin ASCON.
    \item Comparación de dos familias de modelos federables (MLP/RN y CNN-1D) bajo el mismo protocolo HFL, con un Tiny Transformer y un Residual MLP como referencias offline.
    \item Estrategia FOG con agregación intermedia entre gateways Raspberry Pi (líder/peer), que no aparece reportada de forma combinada con ASCON-128 en la literatura revisada.
\end{enumerate}

\begin{longtable}{@{}p{0.18\textwidth} p{0.32\textwidth} p{0.42\textwidth}@{}}
\caption{Relación entre literatura y decisiones de diseño.}
\label{tab:literatura_diseno}\\
\toprule
\textbf{Línea} & \textbf{Referencias} & \textbf{Uso en la tesis} \\
\midrule
\endfirsthead
\toprule
\textbf{Línea} & \textbf{Referencias} & \textbf{Uso en la tesis} \\
\midrule
\endhead
TinyML / IDS IoT & \cite{Tekin2023_OnDeviceEnergy, Amuthadevi2025, Alwaisi2024, HeydariMahmoud2025_TinyML_IDS, Dini2024_EmbeddedIIoT_Review} & Selección de modelos compactos, reporte de huella de memoria, latencia de inferencia en ESP32-S3 y elección de \emph{toolchain} de despliegue. \\
Surveys IoT/IIoT seguridad & \cite{Amanullah2020_DL_BigData_IoT_Security, Ni2024_IIoT_Security_Survey, Bagaa2020_ML_IoT_Framework, Huda2018_DBN_SCADA} & Encuadre del problema, justificación de jerarquía Edge--Fog--Cloud y comparación contra enfoques no edge. \\
FL / HFL para IDS & \cite{McMahan2017FedAvg, Imteaj2022_FL_ResourceConstrainedSurvey, Albanbay2025, Ramadan2025, Hajj2023_CrossLayerFL_IDS, Roy2023_FL_AdaptedModel, Chaurasia2024_FL_ResNet_IIoT, Ficco2024_FL_TinyML_Onboard} & Diseño de FedAvg jerárquico, discusión de datos \textit{non-IID} y sustento del entrenamiento on-board en MCUs/RPi. \\
Robustez/escalabilidad FL en MCUs & \cite{Tran2026_EdgeTrustShard} & Justificación de la jerarquía y delimitación de las amenazas no cubiertas (clientes bizantinos). \\
Criptografía ligera & \cite{NISTSP8002322025, Gusita2025_LWC_Survey, Kaur2025_ASCON_Survey, Mirzaali2025_LWC_Permutations, NISTIR8454_2023} & Integración de ASCON-128, comparación contra baseline sin cifrado, justificación de no usar permutaciones \emph{ad-hoc} y delimitación de canales laterales como trabajo futuro. \\
Datasets IoT-IDS & \cite{Moustafa2020_TONIoT, Neto2023_CICIoT2023, Ferrag2022_EdgeIIoTset, Stratosphere_IoT23, Rani2024_IoT_IDS_Datasets, Raminenei2023_ML_IoT_UNSWNB15} & Selección de las 13 features de flujo y de las tres clases (\texttt{normal}, \texttt{mqtt\_bruteforce}, \texttt{scan\_A}). \\
Pre-agregación / data fusion & \cite{Wang2019_DataFusion_CPSS} & Justificación conceptual del modo FOG (FedAvg de gateways antes del PC). \\
\bottomrule
\end{longtable}

\newpage
